% !TeX root = ../TCSVT_template/manuscript.tex
% ALV-CL rewrite informed by CineDub Sec. 3.3.1:
% distinct constraints -> staged conditioning -> transfer -> ordering rationale.
% Retain one conditioning equation and cite the existing backbone objective.
% Stage II includes speaker conditioning; visual and CTC supervision are inactive.
% The ordering rationale is a design motivation, not evidence of gradient
% conflict, catastrophic forgetting, or superiority to other training orders.

\subsection{Acoustic-to-Linguistic-to-Visual Curriculum Learning}
\label{subsec:progressive_training}
\label{sec:progressive_training}

Linguistic and visual conditioning constrain speech generation in different ways: text specifies ordered content without frame-level timing, whereas lip motion provides temporal cues. An audio-only pretrained generator has a speech prior but lacks an explicit transcript-to-speech mapping. Direct dubbing adaptation therefore introduces linguistic control and audiovisual correspondence within the same optimization stage.

We introduce Acoustic-to-Linguistic-to-Visual Curriculum Learning (ALV-CL), following the staged capability acquisition principle explored in CineDub~\cite{dai2026cinedub}. Let $\mathbf z_{\mathrm{ctx}}$, $\mathbf y$, $\mathbf v$, and $\mathbf s$ denote acoustic context, text tokens, lip-motion features, and the speaker embedding. The three stages progressively expand the conditioning sets:
\begin{equation}
\begin{aligned}
\mathbf c^{\mathrm A} &= (\mathbf z_{\mathrm{ctx}}),\\
\mathbf c^{\mathrm{AL}} &= (\mathbf z_{\mathrm{ctx}},\mathbf y,\mathbf s),\\
\mathbf c^{\mathrm{ALV}} &= (\mathbf z_{\mathrm{ctx}},\mathbf y,\mathbf v,\mathbf s).
\end{aligned}
\label{eq:alv_conditioning}
\end{equation}
Speech remains the generation target across all three stages.

Stage I uses a LibriSpeech-pretrained acoustic generator~\cite{panayotov2015librispeech} that reconstructs masked speech from retained audio context. Stage II transfers its compatible acoustic parameters and trains on paired text and speech from LibriTTS~\cite{zen2019libritts}. Audio--text joint attention and speaker conditioning are optimized through masked flow matching, with visual input and CTC supervision disabled. Stage III initializes the dubbing model from this text-conditioned checkpoint and activates visual interaction and Gaussian-windowed audiovisual attention. All generator pathways are then jointly adapted on synchronized video, transcripts, and speech using the objective in~\eqref{eq:mmdit_training_objective}.

The intermediate stage makes transcript-conditioned synthesis available before visual timing is introduced. Audiovisual adaptation can consequently refine an existing text--speech mapping under lip-motion constraints, with linguistic and acoustic parameters remaining trainable. This ordering separates the initial acquisition of linguistic control from its subsequent audiovisual adaptation, while retaining implicit alignment without externally supplied phoneme durations. The curriculum adds no modules or computational stages at inference.
